\chapter{Extended Modal Operators}
\label{sec:extended-modal}
% ===================================================================

The decoration is, at present, external data attached to the rules.  We
now show that it can be \emph{internalized into the logic itself} by
extending the modal operators to carry the dynamical state of a
transition.

Because Chapter~\ref{sec:weighted} replaced the weight map and the cost
map by a single semiring-valued decoration, the extended operator is
correspondingly simpler than the version this chapter previously carried:
it holds one value, not two.  Where both an amplitude and a charge are
wanted at once, one takes $\Semi$ to be a product semiring --- say
$\CC \times (\Real\cup\{\infty\})^k$ --- and the two coordinates ride
together.  Nothing about the logic changes when they do.

\begin{definition}[Extended Modal Operator]
\label{def:ext-modal}
The formulae of \emph{extended HML} are generated by adding to
Definition~\ref{def:cdHML} the operator
\[
  \langle K,\, \dec,\, \dec^{\dagger},\, \vect{A} \rangle\,\phi
\]
where $K$ is a minimal context as before, $\dec \in R$ is the
\emph{forward decoration} of the transition, $\dec^{\dagger} \in R$ its
\emph{reversal}, and $\vect{A}$ the account available at the surface on
which the step draws.  The satisfaction clause is
\[
  \langle P, \trace, \vect{A} \rangle \;\sat\;
  \langle K, \dec, \dec^{\dagger}, \vect{A}_0 \rangle\,\phi
\]
just when:
\begin{enumerate}[label=(\roman*)]
  \item $K[P]$ evolves in one step via some rule $r$ to $P'$, with redex
        $u$;
  \item the decoration assigns $\dec_r(u) = \dec$, where $\dec_r(u)$ is
        evaluated at the most refined witness $\phi_r(u)$ of
        Definition~\ref{def:refcomplete};
  \item the reversal satisfies $\dec^{\dagger} = \overline{\dec}$ in the
        complex coordinate and $\dec^{\dagger} = \dec$ in the tropical
        coordinates --- a debit reversed is a credit of the same size;
  \item $\vect{A}_0 = \vect{A}$, the account recorded in the formula being
        the one the step actually draws on;
  \item the step is affordable in the tropical coordinates, and the drawer
        holds the capability to the account (Remark~\ref{rmk:ambient});
  \item $\langle P', \trace', \vect{A} - \dec \rangle \sat \phi$, the
        subtraction taken in the tropical coordinates only.
\end{enumerate}
\end{definition}

\section{Internalization of the Decoration}

\begin{proposition}
In extended HML the decoration is recoverable from the logic:
$\dec_r(\phi) = \dec$ where
$\phi = \langle K, \dec, \dec^{\dagger}, \vect{A}\rangle\,\psi$.  The
decorated rewrite rules and the extended modal logic are two presentations
of the same structure.
\end{proposition}

This is a small instance of a pattern that recurs throughout this book,
and it is worth naming here because it will be doing heavy work later.
Structure attached to a calculus from outside can often be pushed back
inside it, and the two erasures --- \emph{forgetting} the decoration and
\emph{internalising} it --- are genuinely different operations with
different costs~\cite{TwoErasures}.  Forgetting loses information.
Internalising loses nothing but has to be paid for in expressiveness of
the host.  When the Pledge insisted that rods and clocks belong inside the
model rather than outside it, this is the operation it was asking for.

\section{Factorization: State vs.\ Rule}

The parameters of the extended operator decompose into two classes:

\begin{center}
\begin{tabular}{ll}
\toprule
\textbf{State-dependent} & \textbf{Rule-dependent} \\
\midrule
$K$ (from MSL applied to the current term) & $\dec$ (from the decoration on the rule) \\
$\vect{A}$ (the account at the surface) & $\dec^{\dagger}$ (its reversal) \\
\bottomrule
\end{tabular}
\end{center}

This is the natural decomposition of the extended logic into
\emph{observables} --- what can be measured now --- and \emph{dynamics} ---
how the system evolves.  Note that with a dynamic decoration
(Definition~\ref{def:dyndec}) the boundary moves: the table travels with
the state, so part of what was rule-dependent becomes state-dependent, and
the logic sees the difference.

\section{Located accounts}

The account $\vect{A}$ was previously described as being either
\emph{global} --- one budget for the whole term --- or \emph{local}, with
each parallel component carrying its own, a communication then debiting
``from a shared pool or by negotiated transfer, depending on the resource
semantics chosen.''  That framing is the source of the defect recorded in
Remark~\ref{rmk:ambient}, and it should be retired.  The global/local
distinction is about \emph{where the arithmetic happens}; the distinction
that matters is about \emph{who is entitled to do it}.

In the rho calculus the two questions have the same answer, because a
location is a name.  An account is a message on a channel; the right to
draw on it is the right to receive on that channel; and since fresh
channels are unforgeable, entitlement is possession rather than adjacency.
The extended operator therefore records the account \emph{at a surface},
and clause~(v) of Definition~\ref{def:ext-modal} asks two independent
questions --- is there enough, and may this drawer draw --- where the
earlier presentation asked only the first.

Two consequences are worth flagging.  Conservation laws become genuinely
local, since each located account is conserved separately and transfers
between them are visible as communications rather than as bookkeeping.
And the spatial and temporal structures cease to be independent
coordinates, for the reason given in Remark~\ref{rmk:spacetime}: a draw
now requires both proximity and availability.  The logic is the natural
place to observe this, since a modality that records a located account is
recording a fact about space and a fact about time in a single operator.

% ===================================================================
